\ifdefined\gkver\else\def\gkver{student}\fi
\documentclass[\gkver]{gaokaozhenti}
\begin{document}

\chapter{2018年4月浙江}

\begin{enumerate}
\item
%% number: 1
%% paperName: 2018年4月浙江选考第 1 题 3 分
%% typeId: 1
%% type: 单选题
%% chapter: 运动和力
%% point: 牛顿第一定律
%% method: 无
%% score: 3
%% degree: 950
%% duplicateId: 0
%% body:
通过理想斜面实验得出“力不是维持物体运动的原因”的科学家是\xzanswer{B}

\fourchoices[answer=B]
{亚里士多德}
{伽利略}
{笛卡尔}
{牛顿}

%% answer: B
%% memo:
\memoanswer{
	伽利略通过理想斜面实验得出，力不是维持物体运动的原因，故 B 正确。
}

\item
%% number: 2
%% paperName: 2018年4月浙江选考第 2 题 3 分
%% typeId: 1
%% type: 单选题
%% chapter: 直线运动
%% point: 路程与位移
%% method: 无
%% score: 3
%% degree: 950
%% duplicateId: 0
%% body:
某驾驶员使用定速巡航，在高速公路上以时速 110 公里行驶了 200 公里。其中“时速 110 公里”、“行驶 200 公里”分别是指\xzanswer{D}

\fourchoices[answer=D]
{速度、位移}
{速度、路程}
{速率、位移}
{速率、路程}

%% answer: D
%% memo:
\memoanswer{
	“时速 110 公里”是指汽车的速度大小，即速率；“行驶 200 公里”指汽车运动轨迹的长度，是路程，D 正确。
}

\item
%% number: 3
%% paperName: 2018年4月浙江选考第 3 题 3 分
%% typeId: 1
%% type: 单选题
%% chapter: 运动和力
%% point: 力学单位制
%% method: 无
%% score: 3
%% degree: 850
%% duplicateId: 0
%% body:
用国际单位制的基本单位表示能量的单位，下列正确的是\xzanswer{A}

\fourchoices[answer=A]
{$\Ukg\cdot\Umq/\Usq$}
{$\Ukgmsq$}
{$\UNm$}
{$\UNcdotm$}

%% answer: A
%% memo:
\memoanswer{
	N 不是基本单位，故 CD 错误；由动能的定义式 $E_k=\frac{1}{2}mv^2$ 可知，$\Ukg\cdot\Umq/\Usq$ 为能量的单位，故 A 正确 B 错误。
}

\item
%% number: 4
%% paperName: 2018年4月浙江选考第 4 题 3 分
%% typeId: 1
%% type: 单选题
%% chapter: 曲线运动
%% point: 线速度角速度
%% method: 无
%% score: 3
%% degree: 750
%% duplicateId: 0
%% body:
A、B 两艘快艇在湖面上做匀速圆周运动（如图），在相同的时间内，它们通过的路程之比是 4:3，运动方向改变的角度之比是 3:2，则它们\xzanswer{A}

\onepicture[w=5.0cm]{figs/04.jpg}

\fourchoices[answer=A]
{线速度大小之比为 4:3}
{角速度大小之比为 3:4}
{圆周运动的半径之比为 2:1}
{向心加速度大小之比为 1:2}

%% answer: A
%% memo:
\memoanswer{
	由线速度的定义 $v=\frac{s}{t}$ 得，A、B 两艘快艇的线速度大小之比为 4:3，故 A 正确；由角速度的定义 $\omega=\frac{\theta}{t}$ 得，A、B 两艘快艇的角速度大小之比为 3:2，故 B 错误；由线速度与角速度的关系 $v=\omega r$ 得，A、B 两艘快艇做圆周运动的半径之比为 8:9，故 C 错误；向心加速度 $a=v\omega$，所以 A、B 两艘快艇做圆周运动的向心加速度大小之比为 2:1，故 D 错误。
}

\item
%% number: 5
%% paperName: 2018年4月浙江选考第 5 题 3 分
%% typeId: 1
%% type: 单选题
%% chapter: 电路
%% point: 电功电功率
%% method: 无
%% score: 3
%% degree: 650
%% duplicateId: 0
%% body:
杭州市正将主干道上的部分高压钠灯换成 LED 灯，已知高压钠灯功率为 $400\UW$，LED 灯功率为 $180\UW$，若更换 4000 盏，则一个月可节约的电能约为\xzanswer{B}

\fourchoices[answer=B]
{$9\times10^{2}\UkWh$}
{$3\times10^{5}\UkWh$}
{$6\times10^{5}\UkWh$}
{$1\times10^{12}\UkWh$}

%% answer: B
%% memo:
\memoanswer{
	每天照明大约 12 小时，一个月内节约的电能大约为
	$E=n(P_{\mathrm{N}}-P_{\mathrm{L}})t=4000\times220\times10^{-3}\times30\times12\UkWh\approx3\times10^{5}\UkWh$，故 B 正确。
}

\item
%% number: 6
%% paperName: 2018年4月浙江选考第 6 题 3 分
%% typeId: 1
%% type: 单选题
%% chapter: 电场
%% point: 电荷守恒定律
%% method: 无
%% score: 3
%% degree: 650
%% duplicateId: 0
%% body:
真空中两个完全相同、带等量同种电荷的金属小球 A 和 B（可视为点电荷），分别固定在两处，它们之间的静电力为 $F$。用一个不带电的同样金属球 C 先后与 A、B 球接触，然后移开球 C，此时 A、B 球间的静电力为\xzanswer{C}

\fourchoices[answer=C]
{$\frac{F}{3}$}
{$\frac{F}{4}$}
{$\frac{3F}{8}$}
{$\frac{F}{2}$}

%% answer: C
%% memo:
\memoanswer{
	设最初 A、B 的电荷量均为 $q$。用金属球 C 先后与 A、B 接触后，A、B 的电荷量分别为 $\frac{q}{2}$、$\frac{3q}{4}$。最初，由库仑定律得 $F=k\frac{q^{2}}{r^{2}}$，接触之后，由库仑定律得 $F'=\frac{3}{8}\times k\frac{q^{2}}{r^{2}}=\frac{3}{8}F$，故 C 正确。
}

\item
%% number: 7
%% paperName: 2018年4月浙江选考第 7 题 3 分
%% typeId: 1
%% type: 单选题
%% chapter: 磁场
%% point: 左手定则
%% method: 无
%% score: 3
%% degree: 850
%% duplicateId: 0
%% body:
处于磁场 $B$ 中的矩形金属线框可绕轴 OO$'$ 转动，当线框中通以电流 $I$ 时，如图所示，此时线框左右两边受安培力 $F$ 的方向正确的是\xzanswer{D}

\onepicture[w=5.2cm]{figs/07.png}

\fourchoices[answer=D, ispicture=true, w=3.4cm]
{figs/07a.png}
{figs/07b.png}
{figs/07c.png}
{figs/07d.png}

%% answer: D
%% memo:
\memoanswer{
	由左手定则可知，线框左边所受安培力 $F$ 竖直向上，右边所受安培力 $F$ 竖直向下，故 D 正确。
}

\item
%% number: 8
%% paperName: 2018年4月浙江选考第 8 题 3 分
%% typeId: 1
%% type: 单选题
%% chapter: 运动和力
%% point: 超重失重
%% method: 无
%% score: 3
%% degree: 650
%% duplicateId: 0
%% body:
如图所示，小芳在体重计上完成下蹲动作。下列 $F-t$ 图像能反映体重计示数随时间变化的是\xzanswer{C}

\onepicture[h=3.4cm]{figs/08.jpg}

\fourchoices[answer=C, ispicture=true, w=3.4cm]
{figs/08a.png}
{figs/08b.png}
{figs/08c.png}
{figs/08d.png}

%% answer: C
%% memo:
\memoanswer{
	下蹲过程，小芳先加速下降，处于失重状态，后减速下降，处于超重状态，故 $F$ 先减小后增大，故 C 正确。
}

\item
%% number: 9
%% paperName: 2018年4月浙江选考第 9 题 3 分
%% typeId: 1
%% type: 单选题
%% chapter: 曲线运动
%% point: 天体运动
%% method: 无
%% score: 3
%% degree: 550
%% duplicateId: 0
%% body:
土星最大的卫星叫“泰坦”（如图），每 16 天绕土星一周，其公转轨道半径为 $1.2\times10^{6}\Ukm$。已知引力常量 $G=6.67\times10^{-11}\UN\cdot\Umq/\Ukg^{2}$，则土星的质量约为\xzanswer{B}

\onepicture[w=3.3cm]{figs/09.jpg}

\fourchoices[answer=B]
{$5\times10^{17}\Ukg$}
{$5\times10^{26}\Ukg$}
{$7\times10^{33}\Ukg$}
{$4\times10^{36}\Ukg$}

%% answer: B
%% memo:
\memoanswer{
	由万有引力提供向心力有 $G\frac{Mm}{r^{2}}=m\left(\frac{2\pi}{T}\right)^{2}r$，得 $M=\frac{4\pi^{2}r^{3}}{GT^{2}}$，其中 $T=16\times24\times3600\Us\approx1.38\times10^{6}\Us$，解得 $M\approx5\times10^{26}\Ukg$，故 B 正确。
}

\item
%% number: 10
%% paperName: 2018年4月浙江选考第 10 题 3 分
%% typeId: 1
%% type: 单选题
%% chapter: 直线运动
%% point: 多段匀变速直线运动
%% method: 无
%% score: 3
%% degree: 450
%% duplicateId: 0
%% body:
如图所示，竖井中的升降机可将地下深处的矿石快速运送到地面。某一竖井的深度为 $104\Um$，升降机运行的最大速度为 $8\Ums$，加速度大小不超过 $1\Umsq$。假定升降机到井口的速度为 0，则将矿石从井底提升到井口的最短时间是\xzanswer{C}

\onepicture[w=4.0cm]{figs/10.jpg}

\fourchoices[answer=C]
{$13\Us$}
{$16\Us$}
{$21\Us$}
{$26\Us$}

%% answer: C
%% memo:
\memoanswer{
	矿石先以最大加速度加速上升，再匀速上升，最后以最大加速度减速上升所用时间最短。在加速阶段，所需时间 $t_{1}=\frac{v}{a}=8\Us$，通过的位移为 $x_{1}=\frac{v^{2}}{2a}=32\Um$。在减速阶段与加速阶段相同，在匀速阶段所需时间为 $t_{2}=\frac{x-2x_{1}}{v}=5\Us$。总时间为 $t=2t_{1}+t_{2}=21\Us$。故选 C。
}

\item
%% number: 11
%% paperName: 2018年4月浙江选考第 11 题 3 分
%% typeId: 1
%% type: 单选题
%% chapter: 电场
%% point: 电场中的图象
%% method: 无
%% score: 3
%% degree: 550
%% duplicateId: 0
%% body:
一带电粒子仅在电场力作用下从 A 点开始以 $-v_{0}$ 做直线运动，其 $v-t$ 图像如图所示。粒子在 $t_{0}$ 时刻运动到 B 点，$3t_{0}$ 时刻运动到 C 点，以下判断正确的是\xzanswer{C}

\onepicture[w=4.4cm]{\includesvg{figs/11}}

\fourchoices[answer=C]
{A、B、C 三点的电势关系为 $\varphi_{B}>\varphi_{A}>\varphi_{C}$}
{A、B、C 三点的场强大小关系为 $E_{C}>E_{B}>E_{A}$}
{粒子从 A 点经 B 点运动到 C 点，电势能先增加后减少}
{粒子从 A 点经 B 点运动到 C 点，电场力先做正功后做负功}

%% answer: C
%% memo:
\memoanswer{
	粒子电性未知，则不能判断电场的方向：若电场方向从 B 指向 A，则 $\varphi_{B}>\varphi_{A}>\varphi_{C}$；若电场方向从 A 指向 B，则 $\varphi_{B}<\varphi_{A}<\varphi_{C}$，故 A 错误；在 $v-t$ 图像中，图线的斜率表示加速度，由图知 $k_{B}>k_{A}>k_{C}$，故 $E_{B}>E_{A}>E_{C}$，故 B 错误；粒子从 A 到 B 过程，电场力做负功，动能减少，电势能增加；从 B 到 C 过程，电场力做正功，动能增加，电势能减少，故 C 正确 D 错误。
}

\item
%% number: 12
%% paperName: 2018年4月浙江选考第 12 题 3 分
%% typeId: 1
%% type: 单选题
%% chapter: 磁场
%% point: 磁场的叠加
%% method: 无
%% score: 3
%% degree: 450
%% duplicateId: 0
%% body:
在城市建设施工中，经常需要确定地下金属管线的位置，如图所示。有一种探测方法是，首先给金属长直管线通上电流，再用可以测量磁场强弱、方向的仪器进行以下操作：①用测量仪在金属管线附近的水平地面上找到磁感应强度最强的某点，记为 a；②在 a 点附近的地面上，找到与 a 点磁感应强度相同的若干点，将这些点连成直线 EF；③在地面上过 a 点垂直于 EF 的直线上，找到磁场方向与地面夹角为 $45^\circ$ 的 b、c 两点，测得 b、c 两点距离为 $L$。由此可确定金属管线\xzanswer{A}

\onepicture[w=4.2cm]{\includesvg{figs/12}}

\fourchoices[answer=A]
{平行于 EF，深度为 $\frac{L}{2}$}
{平行于 EF，深度为 $L$}
{垂直于 EF，深度为 $\frac{L}{2}$}
{垂直于 EF，深度为 $L$}

%% answer: A
%% memo:
\memoanswer{
	用测量仪在金属管线附近的水平地面上找到磁感应强度最强的某点，记为 a，说明 a 点离电流最近；找到与 a 点磁感应强度相同的若干点，将这些点连成直线 EF，故说明这些点均离电流最近，电流应该是平行 EF；画出左侧视图，如图所示：

	\onepicture[w=5.5cm]{\includesvg{figs/12_an}}

	b、c 间距为 $L$，且磁场方向与地面夹角为 $45^\circ$，故深度为 $\frac{L}{2}$。故选 A。
}

\item
%% number: 13
%% paperName: 2018年4月浙江选考第 13 题 3 分
%% typeId: 1
%% type: 单选题
%% chapter: 功和能
%% point: 动能势能
%% method: 无
%% score: 3
%% degree: 450
%% duplicateId: 0
%% body:
如图所示，一根绳的两端分别固定在两座猴山的 A、B 处，A、B 两点水平距离为 $16\Um$，竖直距离为 $2\Um$，A、B 间绳长为 $20\Um$。质量为 $10\Ukg$ 的猴子抓住套在绳子上的滑环从 A 处滑到 B 处。以 A 点所在水平面为参考平面，猴子在滑行过程中重力势能最小值约为（绳处于拉直状态）\xzanswer{B}

\onepicture[w=5.2cm]{figs/13.png}

\fourchoices[answer=B]
{$-1.2\times10^{3}\UJ$}
{$-7.5\times10^{2}\UJ$}
{$-6.0\times10^{2}\UJ$}
{$-2.0\times10^{2}\UJ$}

%% answer: B
%% memo:
\memoanswer{
	画出示意图，如图所示。设 AC 长为 $l_{1}$，BC 长为 $l_{2}$，假设 C 为最低点，此时细绳与水平方向的夹角为 $\theta$，则 $h_{AC}=l_{1}\sin\theta$，$h_{BC}=l_{2}\sin\theta$，$l_{1}+l_{2}=20\Um$，$h_{AC}-h_{BC}=2\Um$，$(l_{1}+l_{2})\cos\theta=16\Um$，得 $\cos\theta=\frac{4}{5}$，$\sin\theta=\frac{3}{5}$，$l_{1}=\frac{35}{3}\Um$，$l_{2}=\frac{25}{3}\Um$，故 $h_{AC}=l_{1}\sin\theta=7\Um$，由重力势能的定义知，猴子的最小重力势能为 $E_{p}=0-mgh_{AC}=-700\UJ$，故 B 正确 ACD 错误。

	\onepicture[w=3.4cm]{figs/13_an.jpg}
}

\item
%% number: 14
%% paperName: 2018年4月浙江选考第 14 题 2 分
%% typeId: 2
%% type: 多选题
%% chapter: 原子和原子核
%% point: 天然放射性
%% method: 无
%% score: 2
%% degree: 750
%% duplicateId: 0
%% body:
下列说法正确的是\xzanswer{BD}

\fourchoices[answer=BD]
{组成原子核的核子越多，原子核越稳定}
{$\ce{^238_92 U}$ 衰变为 $\ce{^222_86 Rn}$ 经过 4 次 $\alpha$ 衰变，2 次 $\beta$ 衰变}
{在 LC 振荡电路中，当电流最大时，线圈两端电势差也最大}
{在电子的单缝衍射实验中，狭缝变窄，电子动量的不确定量变大}

%% answer: BD
%% memo:
\memoanswer{
	比结合能越大，原子核越稳定，比结合能大，组成原子核的核子数并不一定多，故 A 错误；由质量数守恒和电荷数守恒得 $238=222+4x$，$92=86+2x-y$，解得 $x=4$，$y=2$，故 B 正确；在 LC 振荡电路中，当电流最大时，电流的变化率为零，此时线圈两端电势差最小，故 C 错误；由不确定性原理 $\Delta x\cdot\Delta p\geq\frac{h}{4\pi}$ 可知，狭缝变窄，电子动量的不确定量变大，故 D 正确。
}

\item
%% number: 15
%% paperName: 2018年4月浙江选考第 15 题 2 分
%% typeId: 2
%% type: 多选题
%% chapter: 原子和原子核
%% point: 玻尔模型
%% method: 无
%% score: 2
%% degree: 650
%% duplicateId: 0
%% body:
氢原子的能级图如图所示，关于大量氢原子的能级跃迁，下列说法正确的是（可见光的波长范围为 $4.0\times10^{-7}\Um\sim7.6\times10^{-7}\Um$，普朗克常量 $h=6.6\times10^{-34}\UJcdots$，真空中的光速 $c=3.0\times10^{8}\Ums$）\xzanswer{BC}

\onepicture[w=3.6cm]{\includesvg{figs/15}}

\fourchoices[answer=BC]
{氢原子从高能级跃迁到基态时，会辐射 $\gamma$ 射线}
{氢原子处在 $n=4$ 能级，会辐射可见光}
{氢原子从高能级向 $n=3$ 能级跃迁时，辐射的光具有显著的热效应}
{氢原子从高能级向 $n=2$ 能级跃迁时，辐射的光在同一介质中传播速度最小的光子能量为 $1.89\UeV$}

%% answer: BC
%% memo:
\memoanswer{
	氢原子从高能级跃迁到基态时，由 $\Delta E=h\frac{c}{\lambda}$ 可得，辐射光的波长最短为 $\lambda=\frac{6.6\times10^{-34}\times3\times10^{8}}{13.6\times1.6\times10^{-19}}\Um\approx9.1\times10^{-8}\Um$，它的波长比 $\gamma$ 射线长，故不可能辐射 $\gamma$ 射线，故 A 错误；同理可得，氢原子从 $n=4$ 能级向低能级跃迁时，辐射的光的波长最短为 $\lambda_{\min}=\frac{6.6\times10^{-34}\times3\times10^{8}}{12.75\times1.6\times10^{-19}}\Um\approx9.7\times10^{-8}\Um$，最长为 $\lambda_{\max}=\frac{6.6\times10^{-34}\times3\times10^{8}}{0.66\times1.6\times10^{-19}}\Um=1.875\times10^{-6}\Um$，可见光的波长在这范围之内，所以会辐射可见光，故 B 正确；氢原子从高能级向 $n=3$ 能级跃迁时，辐射的光的波长最短为 $\lambda_{\min}=\frac{6.6\times10^{-34}\times3\times10^{8}}{1.51\times1.6\times10^{-19}}\Um\approx8.2\times10^{-7}\Um$，它在红外线区域，所以辐射的光有显著的热效应，故 C 正确；光的波长越长、光相对同一种介质的折射率越小，波速越大，据此可知氢原子从高能级向 $n=2$ 能级跃迁时，辐射的光在同一种介质中传播速度最大的光子能量为 $(3.4-1.51)\UeV=1.89\UeV$，故 D 错误。
}

\item
%% number: 16
%% paperName: 2018年4月浙江选考第 16 题 2 分
%% typeId: 2
%% type: 多选题
%% chapter: 机械波
%% point: 干涉衍射
%% method: 无
%% score: 2
%% degree: 550
%% duplicateId: 0
%% body:
两列频率相同、振幅均为 $A$ 的简谐横波 P、Q 分别沿 $+x$ 和 $-x$ 轴方向在同一介质中传播，两列波的振动方向均沿 $y$ 轴。某时刻两波的波面如图所示，实线表示 P 波的波峰、Q 波的波谷；虚线表示 P 波的波谷、Q 波的波峰。a、b、c 为三个等间距的质点，d 为 b、c 中间的质点。下列判断正确的是\xzanswer{CD}

\onepicture[w=5.6cm]{figs/16.png}

\fourchoices[answer=CD]
{质点 a 的振幅为 $2A$}
{质点 b 始终静止不动}
{图示时刻质点 c 的位移为 0}
{图示时刻质点 d 的振动方向沿 $-y$ 轴}

%% answer: CD
%% memo:
\memoanswer{
	a 处是两列波的波峰与波谷叠加的位置，质点 a 振动减弱，振幅始终为零，故 A 错误；b 处是两列波的振动加强点，质点 b 在不停地振动，最大振幅为 $2A$，故 B 错误；c 处是两列波的波峰与波谷叠加的位置，质点 c 振动减弱，振幅始终为零，因此图示时刻质点 c 的位移为 0，故 C 正确；图示时刻，横波 P 传播到 d 处时，质点 d 沿 $y$ 轴负方向运动，横波 Q 传播到 d 处时，质点 d 沿 $y$ 轴负方向运动，两列波叠加后，质点 d 沿 $y$ 轴负方向运动，故 D 正确。
}

\item
%% number: 17
%% paperName: 2018年4月浙江选考第 17 题 5 分
%% typeId: 4
%% type: 实验题
%% chapter: 曲线运动
%% point: 研究平抛运动实验
%% method: 无
%% score: 5
%% degree: 750
%% duplicateId: 0
%% body:
\begin{enumerate}
	\item
	用图 \ref{20184月浙江17a} 所示装置做“探究功与速度变化的关系”实验时，除了图中已给出的实验器材外，还需要的测量工具有 \tkanswer{C}（填字母）；

	（A）秒表　　（B）天平　　（C）刻度尺　　（D）弹簧测力计

	\item
	用图 \ref{20184月浙江17b} 所示装置做“验证机械能守恒定律”实验时，释放重物前有以下操作，其中正确的是 \tkanswer{AC}（填字母）；

	（A）将打点计时器的两个限位孔调节到同一竖直线上　　（B）手提纸带任意位置　　（C）使重物靠近打点计时器

	\item
	图 \ref{20184月浙江17c} 是小球做平抛运动的频闪照片，其上覆盖了一张透明的方格纸。已知方格纸每小格的边长均为 $0.80\Ucm$。由图可知小球的初速度大小为 \tkanswer{0.70}$\Ums$（结果保留两位有效数字）。
\end{enumerate}
\threepicture[numA=1, numB=2, numC=3, numstyle=arabic, labelA=20184月浙江17a, labelB=20184月浙江17b, labelC=20184月浙江17c, align=right, hA=3.2cm, hB=5.4cm, hC=4.6cm]
{figs/17a.png}
{figs/17b.png}
{figs/17c.png}
%% answer: 
\jdanswer{
\begin{enumerate}
	\item
	C

	\item
	AC

	\item
	0.70（0.66$\sim$0.74）

\end{enumerate}
}
%% memo:
\memoanswer{
	（1）实验中必须要用刻度尺测量纸带上计数点的间距，故 C 正确。

	（2）要使实验误差尽可能小，释放重物前要将打点计时器的两个限位孔调节到同一竖直线上，以尽可能减少摩擦；手要提到纸带的最上端，使纸带竖直，重物要靠近打点计时器释放，尽可能完全利用到纸带，故 AC 正确。

	（3）小球在水平方向做匀速直线运动，有 $x=v_{0}T$；在竖直方向做自由落体运动，有 $\Delta h=gT^{2}$，得 $v_{0}=x\sqrt{\frac{g}{\Delta h}}$，其中 $x=0.028\Um$，$\Delta h=0.016\Um$，解得 $v_{0}=0.70\Ums$。
}

\item
%% number: 18
%% paperName: 2018年4月浙江选考第 18 题 5 分
%% typeId: 4
%% type: 实验题
%% chapter: 电路
%% point: 伏安法测电阻
%% method: 无
%% score: 5
%% degree: 650
%% duplicateId: 0
%% body:
\begin{enumerate}
	\item
	小明用多用电表测量一小段 2B 铅笔芯的电阻 $R_{x}$，正确的操作顺序是 \tkanswer{EDBCA}（填字母）；

	（A）把选择开关旋转到交流电压最高档

	（B）调节欧姆调零旋钮使指针指到欧姆零点

	（C）把红、黑表笔分别接在 $R_{x}$ 两端，然后读数

	（D）把选择开关旋转到合适的档位，将红、黑表笔接触

	（E）把红、黑表笔分别插入多用电表“+”、“$-$”插孔，用螺丝刀调节指针定位螺丝，使指针指 0

	\item
	小明正确操作后，多用电表的指针位置如图所示，则 $R_{x}=$ \tkanswer{2.9}$\UO$。

	\item
	小张认为用多用电表测量小电阻误差太大，采用伏安法测量。

	现有实验器材如下：电源（电动势 $3\UV$，内阻可忽略），电压表（量程 $3\UV$，内阻约为 $3\UkO$），多用电表（$2.5\UmA$ 挡、$25\UmA$ 挡和 $250\UmA$ 挡，对应的内电阻约为 $40\UO$，$4\UO$ 和 $0.4\UO$），滑动变阻器 $R_{P}$（$0\sim10\UO$），定值电阻 $R_{0}$（阻值 $10\UO$），开关及导线若干。

	测量铅笔芯的电阻 $R_{x}$，下列电路图中最合适的是 \tkanswer{B}（填字母），多用电表选择开关应置于 \tkanswer{250}$\UmA$ 挡。
\end{enumerate}
\onepicture[w=6.5cm, align=right]{figs/18a.png}
\onepicture[w=11.0cm, align=right]{figs/18b.png}
%% answer: 
\jdanswer{
\begin{enumerate}
	\item
	EDBCA

	\item
	2.9（2.8$\sim$3.0）

	\item
	B，$250\UmA$

\end{enumerate}
}
%% memo:
\memoanswer{
	（1）多用电表使用前，要进行机械调零；每次换挡后，都要表笔短接进行欧姆调零；测量完毕，将选择开关置于交流电压最高挡。

	（2）根据多用电表的读数规则可得，读数为 $2.9\UO$。

	（3）铅笔芯的电阻较小，因此电流表采用外接法测量误差较小，故 A 错误；实验中，为防止电流表被烧毁，用定值电阻 $R_{0}$ 作为保护电阻，且为了使电压表读数误差小，电压表要跨接在定值电阻与铅笔芯电阻的两端，如选项 B 所示；若连接成选项 D 所示电路，则电源可能会出现短接情况，损毁电源；$I=\frac{U}{R_{0}+R_{x}}\approx230\UmA$，所以选择开关置于 $250\UmA$ 挡。
}

\item
%% number: 19
%% paperName: 2018年4月浙江选考第 19 题 9 分
%% typeId: 6
%% type: 计算题
%% chapter: 运动和力
%% point: 牛顿定律的应用
%% method: 无
%% score: 9
%% degree: 550
%% duplicateId: 0
%% body:
可爱的企鹅喜欢在冰面上游玩，如图所示，有一企鹅在倾角为 $37^\circ$ 的斜面上，先以加速度 $a=0.5\Umsq$ 从冰面底部由静止开始沿直线向上“奔跑”，$t=8\Us$ 时，突然卧倒以肚皮贴着冰面向前滑行，最后退滑到出发点，完成一次游戏（企鹅在滑动过程中姿势保持不变）。若企鹅肚皮与冰面间的动摩擦因数 $\mu=0.25$，已知 $\sin37^\circ=0.6$，$\cos37^\circ=0.8$。求：
\begin{enumerate}
	\item
	企鹅向上“奔跑”的位移大小；
	\item
	企鹅在冰面上滑动的加速度大小；
	\item
	企鹅退滑到出发点的速度大小。（计算结果可用根式表示）
\end{enumerate}
\onepicture[w=5.5cm, align=right]{figs/19.png}
%% answer: 
\jdanswer{
\begin{enumerate}
	\item
	$x=16\Um$

	\item
	$a_{1}=8\Umsq$，$a_{2}=4\Umsq$

	\item
	$v=2\sqrt{34}\Ums\approx11.7\Ums$

\end{enumerate}
}
%% memo:
\memoanswer{
	（1）在企鹅“奔跑”过程中，由运动学公式得 $x=\frac{1}{2}at^{2}$，解得 $x=16\Um$。

	（2）在企鹅上滑过程中，由牛顿第二定律得 $a_{1}=g\sin\theta+\mu g\cos\theta$，解得 $a_{1}=8\Umsq$；在企鹅下滑过程中，由牛顿第二定律得 $a_{2}=g\sin\theta-\mu g\cos\theta$，解得 $a_{2}=4\Umsq$。

	（3）由运动学公式得，企鹅上滑位移为 $x_{1}=\frac{(at)^{2}}{2a_{1}}=1\Um$，退滑到出发点的速度满足 $v^{2}=2a_{2}(x+x_{1})$，联立解得 $v=2\sqrt{34}\Ums\approx11.7\Ums$。
}

\item
%% number: 20
%% paperName: 2018年4月浙江选考第 20 题 12 分
%% typeId: 6
%% type: 计算题
%% chapter: 曲线运动
%% point: 平抛运动
%% method: 极值
%% score: 12
%% degree: 350
%% duplicateId: 0
%% body:
如图所示，一轨道由半径为 $2\Um$ 的四分之一竖直圆弧轨道 AB 和长度可以调节的水平直轨道 BC 在 B 点平滑连接而成。现有一质量为 $0.2\Ukg$ 的小球从 A 点无初速度释放，经过圆弧上的 B 点时，传感器测得轨道所受压力大小为 $3.6\UN$，小球经过 BC 段所受阻力为其重力的 0.2 倍，然后从 C 点水平飞离轨道，落到水平面上的 P 点，P、C 两点间的高度差为 $3.2\Um$。小球运动过程中可以视为质点，且不计空气阻力。
\begin{enumerate}
	\item
	求小球运动至 B 点的速度大小；
	\item
	求小球在圆弧轨道上克服摩擦力所做的功；
	\item
	为使小球落点 P 与 B 点的水平距离最大，求 BC 段的长度；
	\item
	小球落到 P 点后弹起，与地面多次碰撞后静止。假设小球每次碰撞机械能损失 75\%，碰撞前后速度方向与地面的夹角相等。求小球从 C 点飞出后静止所需的时间。
\end{enumerate}
\onepicture[h=5.0cm, align=right]{figs/20.png}
%% answer: 
\jdanswer{
\begin{enumerate}
	\item
	$v_{B}=4\Ums$

	\item
	$W_{f}=2.4\UJ$

	\item
	$L_{BC}=3.36\Um$

	\item
	$t=2.4\Us$

\end{enumerate}
}
%% memo:
\memoanswer{
	（1）小球运动至 B 点，由牛顿第二定律得 $F_{N}-mg=\frac{mv_{B}^{2}}{R}$，解得 $v_{B}=4\Ums$。

	（2）小球由 A 至 B 的过程，由动能定理得 $mgR-W_{f}=\frac{1}{2}mv_{B}^{2}$，解得 $W_{f}=2.4\UJ$。

	（3）小球由 B 至 C 的过程，由动能定理得 $-kmgL_{BC}=\frac{1}{2}mv_{C}^{2}-\frac{1}{2}mv_{B}^{2}$，故 B 至 P 的水平距离为 $L=\frac{v_{B}^{2}-v_{C}^{2}}{2kg}+v_{C}\sqrt{\frac{2h}{g}}=4-\frac{1}{4}v_{C}^{2}+\frac{4}{5}v_{C}$，当 $v_{C}=1.6\Ums$ 时 P 至 B 的水平距离最大，解得 $L_{BC}=3.36\Um$。

	（4）小球从 C 至 P 运动的时间，由运动学公式得 $t_{0}=\sqrt{\frac{2h}{g}}=0.8\Us$；与地面碰撞后，小球的机械能损失 75\%，即动能损失 75\%，即水平与竖直方向的动能均损失 75\%，故竖直方向的速度为碰撞前的一半，因此从 P 点弹起到小球第 1 次落地的时间为 $t_{1}=\frac{1}{2}\times2t_{0}=0.8\Us$，从 P 点弹起到小球第 2 次落地的时间为 $t_{2}=\frac{1}{2}\times t_{1}$，$\cdots\cdots$，从 P 点弹起后，小球第 $n$ 次落地时间为 $t_{n}=\left(\frac{1}{2}\right)^{n-1}\times t_{1}$，故小球最后静止时经历的时间为 $t=t_{0}+t_{1}+t_{2}+\cdots+t_{n}=2.4\Us$。
}

\item
%% number: 21
%% paperName: 2018年4月浙江选考第 21 题 4 分
%% typeId: 4
%% type: 实验题
%% chapter: 光学
%% point: 光的衍射
%% method: 无
%% score: 4
%% degree: 750
%% duplicateId: 0
%% body:
\begin{enumerate}
	\item
	细丝和单缝有相似的衍射图样。在相同条件下，小明用激光束分别垂直照射两种不同直径的细丝 Ⅰ 和细丝 Ⅱ，在光屏上形成的衍射图样如图 \ref{20184月浙江21a} 中 a 和 b 所示。已知细丝 Ⅰ 的直径为 $0.605\Umm$，现用螺旋测微器测量细丝 Ⅱ 的直径，如图 \ref{20184月浙江21b} 所示，细丝 Ⅱ 的直径为 \tkanswer{0.998}$\Umm$。图 \ref{20184月浙江21a} 中的 \tkanswer{a}（填“a”或“b”）是细丝 Ⅱ 的衍射图样。
	\item
	小明在做“用双缝干涉测量光的波长”实验时，尝试用单缝和平面镜做类似实验。单缝和平面镜的放置如图 \ref{20184月浙江21c} 所示，白炽灯发出的光经过滤光片成为波长为 $\lambda$ 的单色光照射单缝，能在光屏上观察到明暗相间的干涉条纹。小明测得单缝与镜面延长线的距离为 $h$，与光屏的距离为 $D$，则条纹间距 $\Delta x=$ \tkanswer{$\frac{D}{2h}\lambda$}。随后小明撤去平面镜，在单缝下方 A 处放置同样的另一单缝，形成双缝结构，则在光屏上 \tkanswer{不能}（填“能”或“不能”）观察到干涉条纹。
\end{enumerate}
\threepicture[numA=1, numB=2, numC=3, numstyle=arabic, labelA=20184月浙江21a, labelB=20184月浙江21b, labelC=20184月浙江21c, align=right, hA=3.4cm, hB=3.4cm, hC=2.4cm]
{figs/21a.png}
{figs/21b.png}
{figs/21c.png}
%% answer: 
\jdanswer{
\begin{enumerate}
	\item
	0.998（0.996$\sim$1.000），a

	\item
	$\frac{D}{2h}\lambda$，不能

\end{enumerate}
}
%% memo:
\memoanswer{
	（1）细丝 Ⅱ 的直径测量值为 $0.5\Umm+49.8\times0.01\Umm=0.998\Umm$。在单缝衍射实验中，缝隙越小，衍射条纹越宽；由题意可知，细丝与单缝有相似的衍射图样，因此有“细丝越细，衍射条纹越宽”的特点，故 a 是细丝 Ⅱ 的衍射图样。

	（2）此实验类似于双缝干涉，只是此时“双缝”之间的距离 $d=2h$，由双缝干涉条纹间距公式 $\Delta x=\frac{l}{d}\lambda$ 可得 $\Delta x=\frac{D}{2h}\lambda$。若撤去平面镜，在单缝下方再放置另一单缝，虽然有两个点光源（两个双缝类似于点光源），但由于缺少单缝不能形成稳定的相干光，故无法观察到干涉条纹。
}

\item
%% number: 22
%% paperName: 2018年4月浙江选考第 22 题 10 分
%% typeId: 6
%% type: 计算题
%% chapter: 磁场
%% point: 带电粒子在磁场中的运动
%% method: 无
%% score: 10
%% degree: 350
%% duplicateId: 0
%% body:
压力波测量仪可将待测压力波转换为电压信号，其原理如图 1 所示。压力波 $p(t)$ 进入弹性盒后，通过与铰链 O 相连的“$\dashv$”型轻杆 L，驱动杆端头 A 处的微型霍尔片在磁场中沿 $x$ 轴方向做微小振动，其位移 $x$ 与压力 $p$ 成正比（$x=\alpha p$，$\alpha>0$）。霍尔片的放大图如图 2 所示，它由长$\times$宽$\times$厚 $=a\times b\times d$、单位体积内自由电子数为 $n$ 的 N 型半导体制成。磁场方向垂直于 $x$ 轴向上，磁感应强度大小为 $B=B_{0}(1-\beta|x|)$，$\beta>0$。无压力波输入时，霍尔片静止在 $x=0$ 处，此时给霍尔片通以沿 C$_{1}$C$_{2}$ 方向的电流 $I$，则在侧面上 D$_{1}$、D$_{2}$ 两点间产生霍尔电压 $U_{0}$。
\begin{enumerate}
	\item
	指出 D$_{1}$、D$_{2}$ 两点哪点电势高；
	\item
	推导出 $U_{0}$ 与 $I$、$B_{0}$ 之间的关系式（提示：电流 $I$ 与自由电子定向移动速率 $v$ 之间关系为 $I=nevbd$，其中 $e$ 为电子电荷量）；
	\item
	弹性盒中输入压力波 $p(t)$，霍尔片中通以相同电流，测得霍尔电压 $U_{H}$ 随时间 $t$ 变化图像如图 3。忽略霍尔片在磁场中运动产生的电动势和阻尼，求压力波的振幅和频率。（结果用 $U_{0}$、$U_{1}$、$t_{0}$、$\alpha$ 及 $\beta$ 表示）
\end{enumerate}
\onepicture[w=9.5cm, align=right]{figs/22a.png}
\onepicture[w=5.0cm, align=right]{figs/22b.png}
%% answer: 
\jdanswer{
\begin{enumerate}
	\item
	$D_{1}$ 点电势高

	\item
	$U_{0}=\frac{1}{ne}\cdot\frac{IB_{0}}{d}$

	\item
	$A=\frac{1}{\alpha\beta}\left(1-\frac{U_{1}}{U_{0}}\right)$，$f=\frac{1}{2t_{0}}$

\end{enumerate}
}
%% memo:
\memoanswer{
	（1）电流中的电子由于受到洛伦兹力的作用，向 $D_{2}$ 侧板运动，故 $D_{1}$ 点电势高。

	（2）当电子不再发生偏转时，电子受到的电场力与洛伦兹力平衡，由平衡条件得 $evB_{0}=eE_{H}$，由电势差定义知 $U_{0}=E_{H}b=\frac{1}{ne}\cdot\frac{IB_{0}}{d}$。

	（3）分析 $U_{H}-t$ 图像，霍尔电压与振幅有关。由（2）知 $U_{0}=\frac{1}{ne}\cdot\frac{IB_{0}}{d}$，则 $U=\frac{1}{ne}\cdot\frac{IB}{d}$，化简得 $\frac{U}{U_{0}}=\frac{B}{B_{0}}=1-\beta|x|=1-\alpha\beta|p|$，即霍尔电压 $U_{H}(t)=U_{0}(1-\alpha\beta|p|)$，解得 $|p|=\frac{1}{\alpha\beta}\left(1-\frac{U}{U_{0}}\right)$。当 $U=U_{1}$ 时，振幅最大，因此压力波的振幅为 $A=\frac{1}{\alpha\beta}\left(1-\frac{U_{1}}{U_{0}}\right)$。由表达式与图像可知周期为 $2t_{0}$，因此压力波的频率为 $f=\frac{1}{2t_{0}}$。
}

\item
%% number: 23
%% paperName: 2018年4月浙江选考第 23 题 10 分
%% typeId: 6
%% type: 计算题
%% chapter: 电磁感应
%% point: 电磁感应受力分析
%% method: 无
%% score: 10
%% degree: 350
%% duplicateId: 0
%% body:
如图所示，在竖直平面内建立 $xOy$ 坐标系，在 $0\leq x\leq0.65\Um$、$y\leq0.40\Um$ 范围内存在一具有理想边界，方向垂直纸面向内的匀强磁场区域。一边长 $l=0.10\Um$、质量 $m=0.02\Ukg$、电阻 $R=0.40\UO$ 的匀质正方形刚性导线框 abcd 处于图示位置，其中心的坐标为（0，$0.65\Um$）。现将线框以初速度 $v_{0}=2.0\Ums$ 水平向右抛出，线框在进入磁场过程中速度保持不变，然后在磁场中运动，最后从磁场右边界离开磁场区域，完成运动全过程。线框在全过程中始终处于 $xOy$ 平面内，其 ab 边与 $x$ 轴保持平行，空气阻力不计。求：
\begin{enumerate}
	\item
	磁感应强度 $B$ 的大小；
	\item
	线框在全过程中产生的焦耳热 $Q$；
	\item
	在全过程中，cb 两端的电势差 $U_{cb}$ 与线框中心位置的 $x$ 坐标的函数关系。
\end{enumerate}
\onepicture[w=5.2cm, align=right]{figs/23.png}
%% answer: 
\jdanswer{
\begin{enumerate}
	\item
	$B=2\UT$

	\item
	$Q=0.0375\UJ$

	\item
	进入磁场前 $x\leq0.4\Um$，$U_{cb}=0$；进入磁场过程 $0.4\Um<x\leq0.5\Um$，$U_{cb}=(4x-1.7)\UV$；在磁场中 $0.5\Um<x\leq0.6\Um$，$U_{cb}=0.4\UV$；出磁场过程 $0.6\Um<x\leq0.7\Um$，$U_{cb}=\frac{1-x}{4}\UV$；线框出磁场后 $x>0.7\Um$，$U_{cb}=0$

\end{enumerate}
}
%% memo:
\memoanswer{
	（1）由几何关系知，线框到达磁场时的竖直位移为 $y=(0.65-0.4-0.05)\Um=0.2\Um$，由运动学公式得，线框竖直方向的速度为 $v_{y}^{2}=2gy$，联立解得 $v_{y}=2\Ums$。由法拉第电磁感应定律和欧姆定律知，线框中的感应电流为 $I=\frac{Blv_{y}}{R}$，线框受力平衡，由平衡条件得 $mg=BIl$，联立解得 $B=2\UT$。

	（2）线框出磁场时，水平方向由动量定理得 $-Bl\Delta q=mv-mv_{0}$，又 $\Delta q=\frac{Bl^{2}}{R}$，对线框运动的全过程，由能量守恒得 $Q=mgl+\frac{1}{2}mv_{0}^{2}-\frac{1}{2}mv^{2}$，联立解得 $Q=0.0375\UJ$。

	（3）线框进入磁场前，即 $x\leq0.4\Um$，线框做平抛运动，且无磁场的存在，则 $U_{cb}=0$；线框进入磁场过程，即 $0.4\Um<x\leq0.5\Um$，cd 段有效切割产生的电动势为 $Bv_{0}v_{y}t$；ab 段切割时，相当于电源，此时 cb 段分到的电压为 $I\frac{R}{4}$，则 $U_{cb}=Bv_{0}v_{y}t-I\frac{R}{4}=(4x-1.7)\UV$；线框在磁场中，即 $0.5\Um<x\leq0.6\Um$，cb 段切割产生的电动势为 $U_{cb}=Blv_{0}=0.4\UV$；线框出磁场过程，即 $0.6\Um<x\leq0.7\Um$，线框切割时水平方向的速度为 $v_{x}=v_{0}-\frac{Bl\Delta q'}{m}=5(1-x)\Ums$，由法拉第电磁感应定律与闭合电路的欧姆定律得 $U_{cb}=\frac{Bv_{x}l}{R}\times\frac{R}{4}=\frac{1-x}{4}\UV$；线框出磁场后，即 $x>0.7\Um$，线框做抛体运动，且无磁场的存在，则 $U_{cb}=0$。
}

\end{enumerate}

\end{document}
